\begin{solapartat}
Al ser el camp magnètic perpendicular al pla de la forca tindrem:
\begin{align*}
\Phi(t=0) &= B\,\textit{Àrea}(t=0) = B\,x_0\,L \ \text{\punts{0,3}} = 0,5\ \mathrm{T}\times 1\,\mathrm{m}\times 2\,\mathrm{m} = 1\,\mathrm{Wb} \ \text{\punts{0,2}}\\
\textit{Àrea}(t) &= L\,x(t) = L\,(x_0 - 0,3\sin(32\,t)) \Rightarrow \ \text{\punts{0,2}}\\
\Phi(t) &= B\,L\,(x_0 - 0,3\sin(32\,t)) = 0,5\times 2\times(1 - 0,3\sin(32\,t))\ \mathrm{Wb} \ \text{\punts{0,3}}
\end{align*}
\end{solapartat}

\begin{solapartat}
\[ \varepsilon = \left|\frac{\mathrm{d}\Phi}{\mathrm{d}t}\right| \ \text{\punts{0,2}} = 0,5\times 2\times 0,3\times 32\,\cos(32\,t) = 9,6\,\cos(32t)\ \mathrm{V} \ \text{\punts{0,3}} \]
El seu valor màxim serà:
\[ \varepsilon_{\textit{màxim}} = 9,6\ \mathrm{V} \ \text{\punts{0,5}} \]
\end{solapartat}
